\par\noindent\textbf{(1)} 界面平行分量 \(k_x=n_1k_0\sin\theta\) 连续，下层法向满足 \(k_{z2}^2=n_2^2k_0^2-k_x^2\)。全反射区 \(k_{z2}=i\kappa\)，其中 \(\kappa=k_0\sqrt{n_1^2\sin^2\theta-n_2^2}>0\)，所以透射场沿法向以 \(e^{-\kappa z}\) 衰减，强度穿透长度为 \(1/(2\kappa)\)。
\par\noindent\textbf{(2)} 对于 \(s\) 偏振，\(r_s=(k_{z1}-k_{z2})/(k_{z1}+k_{z2})=(n_1k_0\cos\theta-i\kappa)/(n_1k_0\cos\theta+i\kappa)\)。其模为1，相位为 \(\varphi_s=-2\arctan[\kappa/(n_1k_0\cos\theta)]\)，反射全能量但可能引入与角度有关的相位延迟。
\par\noindent\textbf{(3)} 窄角谱束的反射相位 \(\varphi_s(k_x)\) 导致驻相位置偏移，Artmann公式 \(D_{\mathrm{GH}}=-d\varphi_s/dk_x\)。令 \(k_1=n_1k_0\)、\(k_2=n_2k_0\)，由 \(\varphi_s=-2\arctan[\sqrt{k_x^2-k_2^2}/\sqrt{k_1^2-k_x^2}]\)，求导得到 \(D_{\mathrm{GH}}=2k_x/[\,\sqrt{k_1^2-k_x^2}\sqrt{k_x^2-k_2^2}\,]\)。
\par\noindent\textbf{(4)} 临界角附近 \(\kappa\to0\)，平面波公式形式上使位移发散，但有限束宽角谱跨越临界角，驻相近似失效并截断该发散。掠入射 \(k_{z1}\to0\) 也需考虑有限波束及界面条件；该位移不是自由空间超光速能流。
